\chapter*{Preface}
Chance seems the opposite of mathematics, and yet it obeys laws as precise as any. This volume of \emph{Order and
Chance} starts with puzzles that a few lines of counting settle, and paradoxes where intuition fails, and it climbs
to the theorems that make randomness predictable: the law of large numbers, the central limit theorem, the
convergence of Markov chains and of martingales, and Brownian motion, the continuous curve that a random walk
becomes when it is seen from far away. It ends with the probabilistic method, where chance proves that objects
exist without ever building them.

Chapters~1--4 live in the discrete world, where sums replace integrals: counting and paradoxes, random variables,
generating functions, and the inequalities that lead to the law of large numbers and to entropy. Chapters~5 and~6
build continuous probability properly and prove the central limit theorem. Chapters~7--11 follow random processes
in time, from random walks and Markov chains to Poisson processes, martingales and Brownian motion. Chapter~12
turns probability into a tool for combinatorics.

The foundation of modern probability is measure theory, the subject of Volume~7. This volume sets up the axioms in
Chapter~1 and builds the integral in Chapter~5, proving its main theorems, but it uses two facts whose proofs belong
to Volume~7: that the uniform probability on $[0,1]$ and infinite products of probabilities exist, and the
Fubini--Tonelli theorem on iterated integrals. Each time they are used, they are labelled \emph{Admitted (proved in
Volume~7)}. Apart from these, every theorem, proposition and lemma is proved; results quoted only for their interest appear as remarks, with a
reference to where they are proved. Exercises are graded \textbf{$\star$} for a warm-up, \textbf{$\star\star$} for a good
afternoon and \textbf{$\star\star\star$} for gems, and every one has a hint and a complete solution at the end of its
chapter, beginning with \emph{How to find it} before the \emph{Write-up}. ``Pause and think'' boxes ask you to guess
before you read on.

The series asks whether pure order and pure chance can be told apart, and why the primes look random yet obey laws. This floor is the load of uncertainty: the central limit theorem, and a Brownian path. The next floor takes up what exact structure remains when a closed form does not. The primes are not yet these coins.
